\section{Method}\label{sec:method}

\subsection{Fixed Method Ladder}

The current experiment evaluates an inherited method ladder rather than a
route-selection policy. Each row changes one major source of information or one
query-generation step. All non-oracle retrieval rows use the same final answer
depth, $k=5$, in the main answer matrix.

\begin{table*}[t]
\centering
\caption{Canonical method ladder. Call counts are conceptual end-to-end LLM
calls. Cached answer replays may log only the final answer call because the
generation and retrieval stages were precomputed under the same method label.}
\label{tab:method_ladder}
\small
\begin{tabularx}{\textwidth}{llcY}
\toprule
Row & Mode & Calls & Purpose \\
\midrule
0 & \method{llm_only} & 1 & Parametric baseline with no retrieval. \\
1 & \method{rag_simple} & 1 & Raw-question retrieval baseline. \\
2 & \method{rag_hyde} & 2 & Question-only hypothetical legal passage, then retrieval and answer. \\
3 & \method{snap_hyre} & 2 & Primary method: snap answer plus HyRE passage in one generation call, then answer from retrieved evidence. \\
4 & \method{golden_passage} & 1 & Oracle diagnostic: inject the labeled gold passage directly. \\
5 & \method{golden_plus_neighbors} & 1 & Oracle-plus-distractor diagnostic: gold passage plus nearby retrieved neighbors. \\
6 & \method{rag_rewrite} & 2 & Non-HyDE generated-query control using a structured query rewrite. \\
\bottomrule
\end{tabularx}
\end{table*}

The oracle rows are not deployable methods. They answer a different question:
if a model receives the labeled evidence, can it use that evidence to select
the right answer? They are therefore diagnostic controls for answer synthesis,
not retrieval baselines.

\subsection{Snap-HyRE Mechanics}\label{sec:method:snaphyre}

\begin{figure*}[t]
  \centering
  \includegraphics[width=0.95\textwidth]{21_snap_hyre_method_flow.png}
  \caption{Operational \hyre{} flow. The first model call produces both a snap
  answer and a retrieval-shaped hypothetical legal reference passage. Retrieval
  uses the passage. The second model call receives the retrieved evidence and
  the original question; the snap answer is logged but not handed in as
  evidence.}
  \label{fig:method_flow}
\end{figure*}

For a formatted legal question $q$, \hyre{} performs:

\begin{enumerate}[leftmargin=*]
\item \textbf{Snap and passage generation.} The model receives the dataset's
direct-answer prompt plus instructions to emit a required final answer line and
a separate \texttt{\#\# Passage} block. The passage must be short, neutral, and
reference-style: it should describe the controlling rule, statutory condition,
legal relationship, or doctrinal test without using answer labels or an
\texttt{Answer:} header.
\item \textbf{Parsing and repair.} The harness parses the first call into
\texttt{snap\_block}, \texttt{snap\_letter}, and \texttt{hyde\_passage}. Under
\method{NO\_SILENT\_FALLBACK=1}, missing snap final lines, missing passage
blocks, answer artifacts inside the passage, parse failures, and overlong
passages are blockers unless a logged same-model format retry repairs them.
\item \textbf{Retrieval.} The retriever uses only the HyRE passage as the query.
The raw snap answer is not evidence. Retrieval uses the same active corpus,
embedding model, hybrid retrieval stack, and cross-encoder reranking path as
the other retrieval rows.
\item \textbf{Final answer.} The final model call receives the retrieved
evidence and the original formatted question. It must emit the dataset-specific
final answer line, such as \texttt{Answer: (C)} or \texttt{Answer: Yes}.
\end{enumerate}

The active harness mode is \method{snap_hyre}. Legacy logs and older docs may
call the same two-call family \method{rag_snap_hyde_2call}; this paper treats
that string as provenance only. The older three-call
\method{rag_snap_hyde} mode is not the current method.

\subsection{Other Rows}

\paragraph{Raw RAG.}
\method{rag_simple} retrieves with the raw question and then answers with the
standard RAG prompt. It is the primary non-oracle baseline because it uses the
same evidence depth and reader prompt as \hyre{} without spending a generated
query call.

\paragraph{HyDE.}
\method{rag_hyde} generates a hypothetical legal reference passage from the
question alone. It isolates whether any generated passage helps, independent of
the snap-conditioned answer commitment used by \hyre{}.

\paragraph{Rewrite.}
\method{rag_rewrite} asks a query-rewriter prompt to emit structured search
queries, then retrieves over that query set and answers from evidence. In
strict mode, malformed rewrite JSON must be repaired by the same model or the
row is rejected rather than silently falling back to raw-question retrieval.

\paragraph{Oracle controls.}
\method{golden_passage} injects the labeled gold passage. \method{golden_plus_neighbors}
injects the gold passage and fills the remaining context slots with nearby
retrieved neighbors. These rows are especially important in law because a
single labeled passage may be under-specified, while additional legal neighbors
may either supply needed context or introduce distracting authority.

\subsection{Caching and Launch Contract}

The evaluation separates generated-query construction, retrieval-id caching,
and answer replay. Generation caches store the HyDE or HyRE passage and parse
metadata. Retrieval caches store deterministic passage ids at max depth
$k=10$, allowing the paper to score Hit@1/5/10 and replay answer rows at
$k=5$ without rerunning embedding search. For large HousingQA rows, hydrated
document caches let answer replay use cached passage ids without reopening the
full 1.8M-document Chroma collection.

Promoted rows must pass the following gates before appearing as paper-facing
claims:

\begin{itemize}[leftmargin=*]
\item \method{NO\_SILENT\_FALLBACK=1}; no provider/method fallback hidden in a
successful row.
\item Exact final answer lines are present and parseable; same-model
format-only retries are logged as caveats.
\item Retrieval caches have unique row keys, no missing indices, no empty
retrieval rows, enough retrieved ids for the requested $k$, and aligned qrels
where gold ids exist.
\item Detail logs pass health scans for errors, missing predictions, parse
failures, fallback markers, long outputs, and near-cap rows.
\item \claimgate names the detail log, the metric values, and any caveats.
\end{itemize}

\subsection{Metrics}

The downstream metric is exact answer accuracy. Multiple-choice tasks use the
last valid answer line, such as \texttt{Answer: (A)} through
\texttt{Answer: (D)} or \texttt{Answer: (E)} depending on the dataset. HousingQA uses
\texttt{Answer: Yes} or \texttt{Answer: No}. Retrieval exposure is Hit@5 and
MRR@5 where the dataset supplies corpus-aligned gold ids. HousingQA also tracks
multiple-gold statute recall internally, but the main table reports any-gold
Hit@5/MRR@5. MASLegalBench has no official per-question gold passage ids, so
retrieval exposure is reported as a same-source-document proxy.

We also report paired McNemar tests when signed in the gate, but we do not
require every signed row to have a significance comparison. Many of the most
important findings are directional comparisons across retrieval and oracle
diagnostics rather than a single pairwise test.
